\documentclass[10pt,a4paper]{amsart}
\usepackage[margin=19mm]{geometry}
\usepackage{amsmath,amssymb,mathtools}
\usepackage[scr=rsfs,cal=euler]{mathalfa}
\usepackage{xfrac}
\usepackage[unicode,hidelinks,bookmarks=false]{hyperref}
\newcommand{\ie}{i.e.\ }
\newcommand{\cf}{cf.\ }
\newcommand{\resp}{respectively\ }
\newcommand{\Subsection}{\subsection}
\providecommand{\nobreakdash}{\nobreak\hbox{-}}
\setlength{\emergencystretch}{2em}
\multlinegap0pt
\usepackage{lmodern,mathtools}
\usepackage{algorithm}
\usepackage{algpseudocode}
\newtheorem{theorem}{Theorem}[section]
\newtheorem{proposition}[theorem]{Proposition}
\newtheorem{lemma}[theorem]{Lemma}
\newtheorem{corollary}[theorem]{Corollary}
\newtheorem{remark}[theorem]{Remark}
\newtheorem{definition}[theorem]{Definition}
\DeclareMathOperator{\disc}{disc}
\newcommand{\R}{\mathbb R}
\newcommand{\E}{\mathbb E}
\newcommand{\Prob}{\mathbb P}
\newcommand{\norm}[1]{\lVert#1\rVert}
\newcommand{\ip}[2]{\langle#1,#2\rangle}
\newcommand{\lstar}{\log^*}
\setlength{\parindent}{1.15em}
\setlength{\parskip}{1pt}
\setlength{\abovedisplayskip}{6pt plus 2pt minus 2pt}
\setlength{\belowdisplayskip}{6pt plus 2pt minus 2pt}
\allowdisplaybreaks
\begin{document}
\title[Improved Algorithms for Beck--Fiala with Bounded Sets]{Improved Algorithms for Beck--Fiala with Bounded Sets}
\author{Dylan J. Altschuler}
\date{}
\maketitle
\noindent\emph{Source and licence.} Improved Algorithms for Beck--Fiala with Bounded Sets. Version arXiv:2609.19714v1 (CC BY 4.0). This material is distributed under the Creative Commons Attribution 4.0 International licence, http://creativecommons.org/licenses/by/4.0/, as recorded in the arXiv OAI licence metadata for this exact version and shown on the abstract page https://arxiv.org/abs/2609.19714v1. Full rights notice: licenses/arxiv-2609.19714-CC-BY-4.0.txt.\par\smallskip
\noindent\textit{Editorial scope.} Counted unit: the original \texttt{theorem} environment \texttt{thm:main} at source lines 71--77 together with its complete proof at lines 146--188; the delivered document keeps source lines 51--193 so that every referenced label and every citation resolves inside it. Native editable source is the paper's own concatenated TeX; the AMS wrapper is editorial formatting only. No publisher class, upstream script or unrelated artwork is redistributed.
\section{Introduction}

For a matrix $A\in\R^{m\times n}$, write
\[ \disc(A)=\min_{\chi\in\{-1,1\}^n}\norm{A\chi}_\infty. \]
If $A$ is the incidence matrix of a set system and each element belongs to at most $d$ sets, the Beck--Fiala conjecture predicted $\disc(A)=O(\sqrt d)$~\cite{BF}. Note that bounds on the maximum row-sum of $A$, which are not assumed in the original Beck-Fiala setting, correspond to bounds on the maximum set size. We consider the task of obtaining efficient (polynomial in the input size) algorithmic discrepancy guarantees for systems with maximum set size of $s$.

\begin{remark}
During the final stages of the writing of this article, Guo, Fang, and Lu used an automatic AI system to prove the stronger Koml\'os conjecture~\cite{GFL}.  Applied to $A/\sqrt d$, this resolves Beck--Fiala with no restriction on the set sizes.  Their proof is existential; the sharp algorithmic problem remains open and is the focus of this article.
\end{remark}

Building on their earlier work~\cite{BJearly}, Bansal and Jiang~\cite{BJ} gave a polynomial-time algorithm for the Beck--Fiala setting that attains discrepancy
\begin{equation}\label{eq:BJ} O\!\left(\min\{d,\sqrt d+(d\log(2n))^{1/3}\}\right), \end{equation}
and hence $O(\sqrt d)$ for $d\gtrsim\log^2n$~\cite[Theorem~5.1]{BJ}. Altschuler and Tikhomirov subsequently improved this to $d\gtrsim_\eta(\log n)(\log\log n)^{2+\eta}$ with an independent approach via an \textit{online} algorithm~\cite{AT}; see also \cite{AT2} which characterizes online discrepancy in the \textit{random} Beck--Fiala setting.  In a complementary direction, if every row has at most $s$ nonzero entries, Harris and Srinivasan give a randomized polynomial-time bound $O(\sqrt{s\log(2d)})$~\cite[Theorem~5.2]{HS}.

We obtain the existential scale algorithmically in bounded-set regimes far below the preceding sparsity thresholds. Define the $j$th smoothed iterated logarithm inductively by
\begin{equation}\label{eq:logs} \ell_0(n)=n, \qquad \ell_{j+1}(n)=\log(e+\ell_j(n)), \end{equation}
and define
\[ \lstar n:=\min\{j\ge0:\ell_j(n)\le e\}. \]
This convention differs from the usual iterated logarithm by at most an additive absolute constant.


\begin{theorem}[Improved algorithmic discrepancy for Beck--Fiala]\label{thm:main}
Let $A\in[-1,1]^{m\times n}$ have at most $d\ge1$ nonzero entries in each column and at most $s\ge1$ in each row.  For every $z\in[-1,1]^n$ and every fixed integer $k\ge1$, there is $\chi\in\{-1,1\}^n$, produced by an efficient randomized algorithm, such that
\begin{equation}\label{eq:fixed} \norm{A(\chi-z)}_\infty \le C_k\min\!\left\{d,\sqrt d+ \bigl[d\{\log(2s)+\ell_{k+1}(n)\}\bigr]^{1/3}\right\}. \end{equation}
There is also a universal constant $C$ for which
\begin{equation}\label{eq:uniform} \norm{A(\chi-z)}_\infty \le C\min\!\left\{d, \bigl[\sqrt d+\{d\log(2s)\}^{1/3}\bigr](1+\lstar n)\right\}. \end{equation}
In particular, taking $z=0$ bounds discrepancy.
\end{theorem}



\begin{remark}[Extension to hereditary discrepancy]
The \textit{hereditary discrepancy} of a matrix is the largest discrepancy of any sub-matrix obtained by considering a subset of the columns. Applying the above theorem to every subset of columns also gives the same bounds for hereditary discrepancy.
\end{remark}



\begin{corollary}[User-friendly formulation]\label{cor:sqrt}
Fix $a>0$ and an integer $j\ge1$. Let $A\in[-1,1]^{m\times n}$ have at most $d\ge1$ nonzero entries per column and at most $\exp(a\sqrt d)$ per row. Then we obtain algorithmic (hereditary) discrepancy of $O_a(\min\{d,\sqrt d(1+\lstar n)\})$. Moreover, under the further assumption that the sparsity is mildly lower bounded, $d\ge\ell_j(n)$, we obtain algorithmic discrepancy of $O_{a,j}(\sqrt d)$.
\end{corollary}
\begin{proof}
Apply~\eqref{eq:uniform} with $s=\exp(a\sqrt d)$, for which $\log(2s)=O_a(\sqrt d)$. If also $d\ge\ell_j(n)$, then $\ell_{j+1}(n)\le\log(e+d)=O(\sqrt d)$, so the second claim follows from~\eqref{eq:fixed} with $k=j$.
\end{proof}

Let us briefly overview the ideas. In~\cite{BJ}, Bansal and Jiang develop a random walk in the cube $[-1,1]^n$. Coordinates are frozen as they approach signs, and linear constraints control the increase of row discrepancies. The main obstacle is to keep enough directions available for the walk to continue making progress toward a signing. The analysis of Bansal and Jiang controls both the coordinate increments and linear combinations of the increments of different rows. The latter control bounds the probability that many rows simultaneously accumulate large discrepancy.

Our modification produces a partial rounding $x$ from an arbitrary starting point $z$, with $\|A(x-z)\|_\infty=O(b)$ and all coordinates outside an exceptional set of columns $B$ rounded to signs. Proposition~\ref{prop:partial} bounds the probability that any specified collection of columns with disjoint row supports all belongs to $B$. The row-size bound then lets us decompose the submatrix obtained by restricting to $B$ into small blocks that can be rounded separately without adding their errors. Repeating this decomposition yields the iterated logarithm bounds. \\

\textbf{AI disclosure.} ChatGPT substantially aided in the proof of Proposition \ref{prop:partial}, specifically in the technical aspects of adapting the analysis in \cite{BJ}.


\section{Bootstrapping scheme}

\begin{definition}\label{def:columngraph}
The \emph{column graph} $G = ([n],E)$ of an $m\times n$ matrix $A$ has a vertex for each column of $A$, and joins two vertices when there is a row that is supported on both corresponding columns. In particular, independent sets of the column graph are collections of columns with disjoint row supports.
\end{definition}


The following decoupling estimate is proved in Section~\ref{sec:joint}; importantly, it imposes no row-size bound.

\begin{proposition}[Joint partial rounding]\label{prop:partial}
There are universal constants $c,C,c_0,d_0>0$ with the following property. Suppose $A\in[-1,1]^{m\times n}$ has absolute column sums at most $d$, $z\in[-1,1]^n$, as well as $d\ge d_0$ and $C\sqrt d\le b\le c_0d$. Then, there is an efficient algorithm computing $x\in[-1,1]^n$ and $B\subseteq[n]$ such that
\begin{equation}\label{eq:partial-error} \norm{A(x-z)}_\infty\le Cb, \qquad x_j\in\{-1,1\} \quad \forall\, j\notin B, \end{equation}
and, for every independent set $R$ in the column graph,
\begin{equation}\label{eq:joint-tail} \Prob(R\subseteq B)\le\exp\{-cb^3|R|/d\}. \end{equation}
\end{proposition}

We now convert the joint tail into a bound on the sizes of the new discrepancy instances that arise in the bootstrapping. This is the sole place that the row-size bounds are needed.

\begin{lemma}[Small subproblems]\label{lem:components}
Let $A$ have $N$ columns and suppose its column graph $G:=G(A)$ has maximum degree at most $\Delta\ge1$. Let $B\subseteq[N]$ be random, and write $G[B]$ for the subgraph induced by $B$. Suppose that, for some $K>0$ and every independent set $R$ of $G$,
\[ \Prob(R\subseteq B)\le e^{-K|R|}. \]
Then, if $K\ge2\log(4\Delta^2)$, there is some constant $C>0$ for which it holds with probability at least $3/4$ that every connected component of $G[B]$ has size at most
\begin{equation}\label{eq:component-size} C(\Delta+1)\left(1+\frac{\log(4N)}K\right). \end{equation}
\end{lemma}

The idea is to associate with each possible connected component of $G[B]$ a single large independent set $R$ of $G$ that has a specific tree structure and must be a subset of $B$. Union bounding over all such trees will give the result.


\begin{proof}
Suppose a connected component $S$ of $G[B]$ has more than $(\Delta+1)(r-1)$ vertices. Starting from any vertex of $S$, select vertices one at a time from $S$, each at $G[B]$-distance two from the previously selected set. If fewer than $r$ vertices have been selected, their closed neighborhoods contain at most $(\Delta+1)(r-1)$ vertices. Hence some vertex of $S$ is outside those neighborhoods. Connectivity of $S$ gives a vertex at distance exactly two from the selected set, so the process can continue until $r$ vertices have been selected.

The selected set $R$ is an independent set. We use $R$ to construct a rooted tree: start from the first vertex added to $R$ and link each subsequent vertex to an earlier vertex at distance two. Ordering the children at each vertex gives a rooted plane tree. The number of such tree shapes with $r$ vertices is the Catalan number $\frac1r\binom{2r-2}{r-1}\le4^{r-1}$~\cite{StanleyEC2}. There are $N$ choices of label for the root, and at most $\Delta^2$ choices for the label of each child vertex. Thus there are at most $N(4\Delta^2)^{r-1}$ possible sets $R$. Each is contained in $B$ with probability at most $e^{-Kr}$, so a union bound gives for every integer $r\ge1$,
\begin{equation}\label{eq:component-probability} \Prob\bigl(G[B]\text{ has a connected component of size } >(\Delta+1)(r-1)\bigr) \le N(4\Delta^2)^{r-1}e^{-Kr}. \end{equation}
In particular, under the stated condition on $K$, the right-hand side is at most $Ne^{-Kr/2}$. Taking $r=\lceil2\log(4N)/K\rceil$ makes this at most $1/4$ and gives \eqref{eq:component-size}.
\end{proof}

Finally, we record the original linear-algebraic bound of Beck and Fiala; while such language was not used explicitly in their original paper, it is immediate from their proof that the result is algorithmic.

\begin{proposition}[Beck--Fiala~\cite{BF}]
Let $F\in \R^{m\times q}$, with each column bounded in $\ell_1$ by $d$. Then, for every $z\in[-1,1]^q$,
\begin{equation}\label{eq:linear} \min_{\chi\in\{-1,1\}^q} \norm{F(\chi-z)}_\infty \le 2d. \end{equation}
Moreover, this is achieved with an efficient algorithm.
\end{proposition}

We are ready to prove the main theorem.


\begin{proof}[Proof of Theorem~\ref{thm:main}]
We repeatedly apply the partial-rounding algorithm, decompose the remaining columns into blocks, and continue rounding each block from its current fractional point. We first explain this decomposition, then bound the block sizes, and finally choose when to stop recursing and complete the rounding.

Consider a current subproblem: a column submatrix $F$ of $A$ with $N$ columns and a starting point $y\in[-1,1]^N$. Initially, $F=A$, $N=n$, and $y=z$. For an admissible parameter $b$, Proposition~\ref{prop:partial} produces $x$ and $B$ with
\[ \norm{F(x-y)}_\infty\le Cb, \qquad x_j\in\{-1,1\}\quad(j\notin B). \]
Keep the signs outside $B$, and remove any already integral coordinates from $B$. If $B$ is empty, this subproblem is complete. Otherwise, let the blocks be the connected components $S$ of $G[B] := G(F)[B]$, the column graph of $F$ restricted to $B$. For each block, we must round $F_S$, the submatrix obtained by restricting $F$ to the columns in $S$, starting from $x_S$.

By construction, no row has support in two different blocks. Consequently, completing these roundings gives
\[ \norm{F(\chi-y)}_\infty \le Cb+\max_S \norm{F_{S}(\chi_S-x_S)}_\infty. \]
Crucially, the errors from different blocks at the same level do not add. Thus each recursive level of the argument adds at most $Cb$ to the discrepancy. Finally, as every block trivially inherits the original bounds $d$ and $s$, we can apply the same procedure again.

We now choose $b$ to ensure that the blocks are small. Take:
\begin{equation}\label{eq:b-choice} b:=C\left(\sqrt d+\{d\log(2ds)\}^{1/3}\right). \end{equation}
Since $(d\log(2d))^{1/3}=O(\sqrt d)$, this satisfies $b=O(\sqrt d+\{d\log(2s)\}^{1/3})$. If $d<d_0$ or $b>c_0d$, the bound $2d$ from \eqref{eq:linear} already proves both conclusions. Otherwise, Proposition~\ref{prop:partial} applies.

The column graph of a subproblem has maximum degree at most $d(s-1)$. With $\Delta=\max\{1,d(s-1)\}$, our choice of $b$ ensures
\[ K:=cb^3/d\ge2\log(4\Delta^2). \]
Lemma~\ref{lem:components} therefore guarantees, with probability at least $3/4$, that every resulting block has at most
\begin{equation}\label{eq:g} g(N):=\left\lceil Cds\log(eN)\right\rceil \end{equation}
columns. If this test fails, discard the output and rerun the partial-rounding algorithm on the same subproblem $F$, from the same starting point $y$, with fresh randomness.

To track the recursion, let $N_0=n$ and $N_{i+1}=g(N_i)$. After $r$ levels of recursion, every unfinished block has at most $N_r$ columns, and the accumulated discrepancy in each row is at most $Crb$. The following simple estimate will handle both conclusions. Choose $a=Cds\ge8$ large enough that $g(x)\le a(1+\log x)$. For $u\ge1$,
\[ g(a^2u) \le a(1+2\log a+\log u) \le a^2\log(e+u). \]
Here we used $1+2\log a\le a-1$. By induction and \eqref{eq:logs},
\[ N_r\le a^2\ell_r(n)\qquad(r\ge0). \]

It remains to complete the rounding of these blocks. For a nonempty block with $N$ columns, apply Proposition~\ref{prop:partial} with $b_N=C(\sqrt d+\{d\log(2N)\}^{1/3})$. Finally, applying \eqref{eq:joint-tail} to singletons and using a union bound gives
\[ \Prob(B\ne\varnothing) \le N\exp(-cb_N^3/d)\le\frac14. \]
Repeat this final application of Proposition~\ref{prop:partial}, each time from the same initial point of the block and with fresh randomness, until $B=\varnothing$. The accepted output is a complete rounding. Of course, if $b_N>c_0d$ then \eqref{eq:linear} can be used instead. Including the option of always using \eqref{eq:linear}, we can therefore finish any block with additional error at most
\begin{equation}\label{eq:terminal} T(N):=C\min\left\{d,\sqrt d+\{d\log(2N)\}^{1/3}\right\}. \end{equation}
Thus, after $r$ recursive levels followed by these final roundings,
\begin{equation}\label{eq:recurrence} \norm{A(\chi-z)}_\infty \le Crb+T(N_r). \end{equation}

For the fixed-depth bound, take $r=k$. Our size estimate gives
\[ \log(2N_k) \le\log(2a^2)+\log\ell_k(n) \le C\{\log(2ds)+\ell_{k+1}(n)\}. \]
Substituting into \eqref{eq:recurrence}, absorbing $(d\log(2d))^{1/3}$ into $O(\sqrt d)$ proves \eqref{eq:fixed}.

For the uniform bound, take $r=\lstar n$. By definition, $\ell_r(n)\le e$, so every remaining block has at most $ea^2$ columns. Since $T(ea^2)=O(b)$, \eqref{eq:recurrence} gives
\[ \norm{A(\chi-z)}_\infty\le Cb(1+\lstar n). \]
Absorbing the $\log d$ term as above and again comparing with \eqref{eq:linear} proves \eqref{eq:uniform}.

Finally, each acceptance test succeeds with probability at least $3/4$, so retries have constant expected overhead. At each depth the subproblems have disjoint column sets, and there are at most $n$ nonempty subproblems. Thus the algorithm has polynomial expected running time. The geometric tails of the retry counts also give polynomial running time with high probability.
\end{proof}




\section{The modified algorithm}\label{sec:joint}


\begin{thebibliography}{99}
\bibitem[AT]{AT} D.~J. Altschuler and K.~Tikhomirov. \newblock Online Beck--Fiala down to logarithmic sparsity. \newblock arXiv:2607.14238, 2026. \newblock \url{https://arxiv.org/abs/2607.14238}.
\bibitem[AT2]{AT2} D.~J. Altschuler and K.~Tikhomirov. \newblock The threshold for online balancing of i.i.d. binary vectors. \newblock arXiv:2609.14975, 2026. \newblock \url{https://arxiv.org/abs/2609.14975}.
\bibitem[BF]{BF} J.~Beck and T.~Fiala. \newblock ``Integer-making'' theorems. \newblock \emph{Discrete Applied Mathematics}, 3(1):1--8, 1981.
\bibitem[BJ]{BJ} N.~Bansal and H.~Jiang. \newblock Beck--Fiala and Koml\'os bounds beyond Banaszczyk. \newblock In \emph{Proceedings of the 58th Annual ACM Symposium on Theory of Computing}, pages 432--442, 2026. \newblock \url{https://arxiv.org/abs/2508.03961}.
\bibitem[BJearly]{BJearly} N.~Bansal and H.~Jiang. \newblock An improved bound for the Beck--Fiala conjecture. \newblock In \emph{2025 IEEE 66th Annual Symposium on Foundations of Computer Science (FOCS)}, pages 220--232, 2025. \newblock \url{https://arxiv.org/abs/2508.01937}.
\bibitem[GFL]{GFL} S.~Guo, E.~X. Fang, and J.~Lu. \newblock Vector balancing via directional total variation. \newblock arXiv:2609.11189, 2026. \newblock \url{https://arxiv.org/abs/2609.11189}.
\bibitem[HS]{HS} D.~G. Harris and A.~Srinivasan. \newblock The Moser--Tardos framework with partial resampling. \newblock \emph{Journal of the ACM}, 66(5):36:1--36:45, 2019.
\bibitem[StanleyEC2]{StanleyEC2} R.~P. Stanley. \emph{Enumerative Combinatorics, Volume~2}. Cambridge University Press, 1999.
\end{thebibliography}
\end{document}
